% TOP_PROBLEM: {"id": 3432, "title": "Closing Lemma for Diffeomorphism (Dynamical Systems)", "category_id": 7, "category": {"id": 7, "name": "topology", "display_name": "Topology", "description": "Properties preserved under continuous deformations.", "slug": "topology", "order_index": 7, "created_at": "2026-07-31T15:26:25.671Z"}, "status": "open", "problem_number": "OPG-48767", "set_id": 12}
% FIRST_POSTED: 2026-10-01
% ENTRY_KIND: statement_only
% UnsolvedMath ID 3432; source catalogue code OPG-48767
% !TeX program = lualatex
% Definition and sources only; this file is not an LLM solution attempt.
\documentclass[11pt,a4paper]{article}
\usepackage[margin=25mm]{geometry}
\usepackage{fontspec}
\setmainfont{lmroman10-regular.otf}[BoldFont=lmroman10-bold.otf,
 ItalicFont=lmroman10-italic.otf,BoldItalicFont=lmroman10-bolditalic.otf]
\usepackage{amsmath,amssymb,mathtools}
\usepackage{unicode-math}
\setmathfont{latinmodern-math.otf}
\AtBeginDocument{\let\setminus\smallsetminus}
\usepackage{xurl,hyperref}
\hypersetup{colorlinks=true,urlcolor=blue}
\setcounter{secnumdepth}{0}
\setlength{\parindent}{0pt}
\setlength{\parskip}{5pt}
\setlength{\emergencystretch}{3em}
\begin{document}
\section*{Closing Lemma for Diffeomorphism (Dynamical Systems)}\label{3432}
{\small UnsolvedMath ID 3432 · OPG-48767 · Topology · OpenGarden}
\subsection{Definitions and mathematical statement}
Fix an integer $r\ge2$ and a compact smooth manifold $M$ without
boundary. Let $\operatorname{Diff}^r(M)$ be the space of bijective
$C^r$ maps with $C^r$ inverses, equipped with the $C^r$ topology:
maps and their derivatives through order $r$ must be uniformly
close in charts. A point $p$ is nonwandering for $f$ if every open
neighborhood $U$ of $p$ satisfies
\[
                       f^n(U)\cap U\ne\varnothing
                       \quad\text{for some integer }n\ge1.
\]
Write $\Omega(f)$ for this set. The closing conjecture asks whether
\[
 \forall f\ \forall p\in\Omega(f)\ \forall C^r\text{ neighborhoods }V\ni f,
 \quad\exists g\in V\ \exists m\ge1\quad g^m(p)=p.
\]
The source's undefined symbol $\omega_f$ is interpreted as the
nonwandering set in the standard closing-lemma formulation. The point
$p$ itself is to become periodic, and the period is not prescribed.
There is also a $C^\infty$ version using its usual smooth topology.

For flows, the analogous question replaces $f$ by a $C^r$ vector
field, uses its flow to define nonwandering points, and asks for
arbitrarily small $C^r$ perturbations creating a periodic orbit
through a given nonsingular nonwandering point. Equilibria must be
treated separately. These higher-regularity targets retain the
specified perturbation topology.
\subsection{Short English statement}
Can every nonwandering point be made periodic by an arbitrarily small perturbation controlling derivatives through any prescribed higher order?
\subsection{Sources}
\begin{itemize}
\item[S1] ULAM.AI and the UnsolvedMath Contributors, supplied problems.json, record 3432 (OPG-48767), OpenGarden collection. Catalogue identity and source transcription; CC BY 4.0.
\url{https://huggingface.co/datasets/ulamai/UnsolvedMath}.
\item[S2] Open Problem Garden, Closing Lemma for Diffeomorphism (Dynamical Systems). Original problem and discussion; the mathematical formulation below is expanded from the supplied source record.
\url{http://www.openproblemgarden.org/op/closing_lemma_for_diffeomorphism_dynamical_systems}.
\end{itemize}
\end{document}
